%%%%%
%
\section{Correction de discontinuit�s}

\itemstyle{Op\'erateur~:} \operfmt{Correc\_discont}
\index{Correc\_discont}

%%%%%OBJET%%%%
\itemstyle{Objet~:}
D�tecte les discontinuit�s dans les s�ries temporelles.
On consid�re qu'un point provoque une discontinuit� si :\\
1) il est en dehors de bornes plausibles, ou alors.\\
2) l'�cart avec le pr�c�dent point correct est sup�rieur (en
valeur absolue) � un seuil.

Les points ainsi d�tect�s sont (au choix) remplac� par
une valeur sp�ciale ("valeur manquante"), ou par une
interpolation.

%%%%%ARGUMENTS%%%%
\bigskip
%%%%%ENTREES%%%%
\begin{tabarg}
\multicolumn{3}{l}{\titretab {Entr\'ees}}\\
\hline
\$Arg1 & Vecteur, {\tt float} & blabla ()\\
\$Arg2 & Coeff., {\tt float} & blabla ()\\
\$Arg3 & Coeff., {\tt float} & blabla ()\\
\$Arg4 & Coeff., {\tt float} & blabla ()\\
\$Arg5 & Coeff., {\tt Int_32         } & blabla ()\\
\$Arg6 & Coeff., {\tt float} & blabla ()\\
\end{tabarg}

\medskip
%%%%%SORTIES%%%%
\begin{tabarg}
\multicolumn{3}{l}{\titretab {Sorties}}\\
\hline
\$Arg7 & Vecteur, {\tt float} & blabla ($Arg1.unit)\\
\end{tabarg}


%%%%%FORMULE%%%%
\itemstyle{Formule}

%%Correc_discont
\[
Resultat = a-b. \Zeta
\]



%%%%%COMMENTAIRE%%%%
\itemstyle{Commentaire~:}
Cette fonction n'est utile que si il y a une diff�rence
d'amplitude marqu�e entre les sauts normaux ("petits") et
les sauts anormaux ("grands").

%%%%%AUTEUR%%%%
\itemstyle{Auteur~:}
CNRM/GMEI/TRAMM

%%%%%REFERENCE%%%%
\itemstyle{R\'ef\'erences~:}

